% ============================================================
% Main_pages/abstract_heb.tex -- Hebrew abstract
% (typeset inside the hebrew environment by main.tex)
% ============================================================
\section*{תקציר}

מערכות יצירה מועשרת-אחזור (\LR{Retrieval-Augmented Generation, RAG}) מעגנות את תשובות מודל השפה (\LR{LLM}) במסמכים שנשלפים ממאגר ידע חיצוני, ולכן איכות התשובה חסומה מלמעלה באיכות האחזור: אם המסמך הרלוונטי לא אוחזר, המודל אינו יכול להשתמש בו. עומדים לרשותנו שני סוגים של אות אחזור. מאחזרים צפופים מדרגים מסמכים לפי הדמיון בין הייצוג הווקטורי (\LR{embedding}) שלהם לבין ייצוג השאילתה. הם זולים, אך בתחומים מתמחים הם מחטיאים את מסמך המטרה בחלק לא זניח מהשאילתות. שופטי \LR{LLM} מעריכים רלוונטיות בדיוק רב בהרבה, אך כל קריאה אליהם יקרה, והציונים שלהם גסים: מסמכים רבים מקבלים את אותו ציון מרבי.

במחקר זה אנו מנסחים אחזור בעזרת שופט \LR{LLM} כבעיית \emph{חיפוש סדרתי תחת תקציב}. המרחק במרחב הייצוגים משמש פונקציית קירוב זולה, והשופט משמש אורקל יקר. בכל צעד מדיניות החיפוש בוחרת איזה מסמך לשלוח לשופט, במטרה להגיע למסמך המטרה במספר קטן ככל האפשר של קריאות. במחקר המקדים בדקנו 196 שאלות מ-\LR{PubMedQA} מול מאגר משותף של 5,035 מסמכים שנשלפו מאינדקס צפוף של \LR{PubMed} המלא, וכל כמיליון זוגות השאלה--מסמך דורגו על ידי \LR{\mbox{Gemini 2.5 Flash-Lite}}. המחקר המקדים ביסס שלוש עובדות. (א) ציון השופט הוא קירוב אמין אך גס לקרבה למסמך המטרה: המתאם עם המרחק שלילי בכל השאילתות ו-\LR{ROC-AUC} שווה ל-0.997, אך בחציון 18 מסמכים לשאילתה חולקים את הציון המרבי. (ב) השאילתות הן נקודות חריגות גיאומטרית ביחס למסמכים. (ג) מרכז הכובד (צנטרואיד) של המסמכים שהשופט אישר קרוב למסמך המטרה ב-44\% יותר מהשאילתה, אך שימוש סטטי בו משפר את דיוק המקום הראשון בנקודת אחוז אחת לכל היותר (מ-70.4\%~ל-71.4\%), והטיה נאיבית ומקוונת של החיפוש לכיוונו מייקרת את החיפוש.

על בסיס ממצאים אלה נפתח ונבחן: (א) מדיניות חיפוש מונחית שופט, שמופעלת רק כאשר סביר שהמאחזר נכשל, משתמשת במשוב דוחה ממסמכים רלוונטיים שאינם המטרה, וממקמת את המטרה בשיטת טרילטרציה מציוני השופט; (ב) שיפוט ברזולוציה גבוהה יותר, שמפרק שוויונות בציון; (ג) אפיון וחיזוי מוקדם של כישלונות אחזור. המדיניות תושווה באמצעות עקומות עלות--דיוק לאחזור צפוף, לדירוג מחדש בעזרת \LR{cross-encoder} ובעזרת \LR{LLM}, למשוב רלוונטיות מדומה ול-\LR{HyDE}, על מאגרים ביו-רפואיים ועל מאגרים מתחום כללי.

\vspace{6pt}
\noindent\textbf{מילות מפתח:} יצירה מועשרת-אחזור, אחזור צפוף, גיאומטריה של מרחב הייצוגים, \LR{LLM} כשופט, חיפוש תחת תקציב, מענה על שאלות ביו-רפואיות.
